% Generated draft: requires mathematical and source review.
\documentclass[11pt]{article}
\usepackage[margin=25mm]{geometry}
\usepackage{fontspec}
\setmainfont{Latin Modern Roman}
\usepackage{amsmath,amssymb,mathtools,unicode-math}
\setmathfont{Latin Modern Math}
\usepackage{enumitem,needspace,xurl,hyperref}
\setlength{\parindent}{0pt}
\setlength{\parskip}{5pt}
\setlength{\emergencystretch}{3em}
\newcommand{\N}{\mathbb{N}}
\newcommand{\Z}{\mathbb{Z}}
\newcommand{\Q}{\mathbb{Q}}
\newcommand{\R}{\mathbb{R}}
\newcommand{\C}{\mathbb{C}}
\newcommand{\F}{\mathbb{F}}
\newcommand{\A}{\mathbb{A}}
\newcommand{\PP}{\mathbb{P}}
\newcommand{\E}{\mathbb{E}}
\newcommand{\eps}{\varepsilon}
\newcommand{\dd}{\,\mathrm{d}}
\begin{document}
\section*{UnsolvedMath ID: 700658}
\textbf{Source problem code:} catalogue entry
\subsection{Definitions and mathematical statement}
Let \(R\) and \(S\) be associative rings with identity, with \(2\cdot 1_S\) invertible in \(S\), and let \(n\geq 3\) and \(m\geq 2\) be integers. No commutativity assumption is imposed. Write \(GL_n(R)\) for the group of invertible \(n\times n\) matrices over \(R\). For \(i\neq j\) and \(r\in R\), put
\[
t_{ij}(r)=I_n+rE_{ij},\qquad
E_n(R)=\langle t_{ij}(r):i\neq j,\ r\in R\rangle,
\]
where \(E_{ij}\) is the matrix unit, and put
\[
a_{ij}=t_{ij}(1)t_{ji}(-1)t_{ij}(1).
\]
For a two-sided ideal \(I\) of \(R\), the source also introduces the natural reduction homomorphism \(\Lambda_I:GL_n(R)\to GL_n(R/I)\); no particular ideal or condition involving this homomorphism is imposed in the requested problem.

Let
\[
q:GL_n(R)\longrightarrow GL_n(R)/Z(GL_n(R))
\]
be the quotient map, write \(\overline{x}=q(x)\), and, for \(G\leq GL_n(R)\), write \(PG=q(G)\). In particular, \(PE_n(R)\) denotes the image of \(E_n(R)\) in this quotient by the centre of \(GL_n(R)\).

The source uses \(GL(W)=GL_m(S)\), with \(W\) a free \(S\)-module of rank \(m\). It calls a homomorphism \(\Lambda:E_n(R)\to GL(W)\) \emph{standard} when there are \(S\)-modules \(P,Q\), a decomposition
\[
S^m\cong P^{\oplus n}\oplus Q,
\]
a unital ring homomorphism \(\delta:R\to\operatorname{End}P\), its entrywise induced homomorphism
\[
\delta^*:GL_n(R)\longrightarrow GL_n(\operatorname{End}P),
\]
a central idempotent \(f\) of \(\delta(R)\), an antiisomorphism \(\nu\) of \(\delta(R)\), an isomorphism \(g:W\to S^m\), and an embedding
\[
\tau:GL_n(\operatorname{End}P)\longrightarrow GL(gW)
\]
such that, for every \(x\in E_n(R)\),
\[
\Lambda(x)=g^{-1}\tau\!\left(
\delta^*(x)f+
\left({}^{t}\delta^*(x)^{\nu}\right)^{-1}(1-f)
\right)g.
\]
Here \({}^{t}\) denotes transposition, and multiplication by \(f\) and \(1-f\) in the matrix expression uses the corresponding scalar diagonal matrices. The formula is reproduced with the source's notation: the codomain and precise matrix action of \(\nu\), the handedness of the modules, and the compatibility required of \(\tau\) with the displayed module decomposition are not fully specified in the supplied record. These qualifications remain unresolved.

For \(i\neq j\), let
\[
a^*_{ij}=t_{ij}(1_{\operatorname{End}P})
             t_{ji}(-1_{\operatorname{End}P})
             t_{ij}(1_{\operatorname{End}P})
\in GL_n(\operatorname{End}P).
\]
This gives the signed matrix meant by the source's replacement of the scalar entries of \(a_{ij}\) by the corresponding entries over \(\operatorname{End}P\).

The source states the following sufficient criterion: a homomorphism
\[
\Lambda_0:PE_n(R)\longrightarrow GL(W)=GL_m(S)
\]
is induced by a standard homomorphism if there are common choices of \(g\) and \(\tau\), with module data as in the preceding setup, such that
\[
\Lambda_0(\overline{a}_{ij})
 =g^{-1}\tau(a^*_{ij})g
 \qquad\text{for every }i\neq j.
\tag{C}
\]
Here being induced means that \(\Lambda_0\circ(q|_{E_n(R)})\) is standard. This criterion is supplied background, not an independently verified assertion in this draft.

\textbf{Problem (V. M. Petechuk, 1986).} Describe the form, at least in particular cases, of homomorphisms \(\Lambda_0:PE_n(R)\to GL_m(S)\) for which condition \((C)\) fails. Thus no single admissible choice of the module data, \(g\), and \(\tau\) makes the displayed equalities hold simultaneously for all \(i\neq j\). The source does not specify which particular cases should be treated or require existence for every allowed \(R,S,n,m\). Moreover, failure of this sufficient criterion is the requested condition; the supplied statement does not assert that it is equivalent to failure of standardness. A fully intrinsic formulation of the admissible choices in \((C)\) requires clarification of the source's module and embedding conventions.
\subsection{Short English statement}
For associative unital rings \(R,S\), with \(2\) invertible in \(S\), and integers \(n\geq3\), \(m\geq2\), describe, at least in particular cases, homomorphisms \(PE_n(R)\to GL_m(S)\) whose images of all \(\overline{a}_{ij}\) cannot simultaneously have the prescribed form under one common choice of \(g\) and \(\tau\). The source gives that form as a sufficient criterion for being induced by a standard homomorphism. Some conventions in its definition of standardness remain unspecified. The registry labels the problem open and records an unverified AI solution claim; neither assessment is independently verified here.
\subsection{Sources}
\begin{itemize}
\item Local UnsolvedMath registry, permanent ID \(700658\), exact source problem code \texttt{catalogue entry}, title \emph{Kourovka Notebook Problem 10.43}. The supplied record attributes the problem to V. M. Petechuk, proposed in 1986.
\item \emph{The Kourovka Notebook (Unsolved Problems in Group Theory)}, original problem 10.43, 10th issue (1986); supplied reference edition: 21st issue (2026), \texttt{arXiv:1401.0300v48}. The registry's collection description identifies E. I. Khukhro and V. D. Mazurov as editors. This is the supplied authoritative source attribution; the publication has not been independently consulted for this draft.
\item Supplied discussion and status-review note: \emph{[Unverified AI claim]} such homomorphisms are found, attributed to Achyuth Jayadevan, \emph{Preprint}, 2026, \url{https://zenodo.org/records/22852216}. The record describes this as an unverified AI solution claim and says the problem is marked with \(\#\) in the October 2026 version of the 21st issue, \texttt{arXiv:1401.0300v48}. No validity or resolution is asserted here.
\item The supplied status review is dated \texttt{2026-10-16}, assigns canonical status \texttt{open}, and records evidence \texttt{unrefereed\_claim}. This assessment is supplied dated curation, not independently verified literature or source verification.
\end{itemize}
\end{document}
